% ==============================================================================
% MÓDULO: INTRODUCCIÓN
% Archivo: subcapitulos/sec_00_introduccion.tex
% ==============================================================================
\section*{INTRODUCCIÓN}
\label{sec:intro}
\addcontentsline{toc}{section}{INTRODUCCIÓN}

La calidad de atención y el acceso efectivo a los servicios de salud en el primer nivel constituyen pilares insoslayables para la consolidación de la cobertura sanitaria universal y la materialización del derecho humano a la salud. De acuerdo con las directrices conjuntas de la Organización Mundial de la Salud (OMS), la OCDE y el Banco Mundial \cite{who2018}, brindar servicios oportunos, seguros, continuos y centrados en la persona es un imperativo ético indispensable para mitigar las inequidades estructurales en salud. En el contexto de la atención primaria, la percepción que el usuario construye sobre el cuidado que se le prodiga no se limita a un mero juicio de complacencia utilitaria; representa una valoración profunda e intersubjetiva que condiciona la adherencia a los tratamientos, la asistencia oportuna a los controles preventivos y el nivel de confianza depositado en las instituciones públicas del Estado.

En el Perú, la Política Nacional de Calidad en Salud del Ministerio de Salud (MINSA) \cite{minsacalidad2021} reconoce que la persistencia de barreras geográficas, administrativas, culturales y de trato interpersonal en los establecimientos de primer nivel (I-1 a I-4) debilita la capacidad de respuesta del sistema sanitario, provocando la despersonalización del acto de cuidado y prolongadas demoras que impactan negativamente en las familias de menores ingresos. Bajo el marco normativo de la cátedra de Proyecto de Investigación en Salud (EN 486) de la Escuela Profesional de Enfermería de la Universidad Nacional de San Cristóbal de Huamanga (UNSCH), impartida por el Dr. Manglio Aguirre Andrade, la formulación de proyectos en salud exige una articulación analítica rigurosa mediante la técnica del embudo en tres niveles (global, nacional y regional/local) y la fundamentación explícita de sus dimensiones justificatorias:

\paragraph{1. Técnica del embudo contextual y territorial:}
\begin{itemize}
    \item \textbf{Nivel Global e Internacional:} A escala mundial, los sistemas de salud enfrentan una crisis de deshumanización derivada de la sobrecarga burocrática y el predominio tecnicista del acto médico, desatendiendo la vivencia afectiva del paciente y los determinantes socioculturales de la salud \cite{who2018}.
    \item \textbf{Nivel Nacional (Perú):} A nivel nacional, el 78\% de las atenciones asistenciales recaen en el primer nivel; no obstante, coexisten cuellos de botella en la asignación de citas, colas nocturnas y precariedad en el trato interpersonal, lo cual motivó al Estado peruano a priorizar la investigación en sistemas sanitarios a través de la \textbf{Línea 10 de Investigación en Salud al 2030}: \textit{``Sistemas y servicios de salud, acceso y cobertura universal''} (aprobada mediante Resolución Ministerial N° 424-2025/MINSA) \cite{minsainv2025}.
    \item \textbf{Nivel Regional y Local (Ayacucho -- Huamanga):} En el plano regional, la investigación se enmarca y responde de manera vinculante a la \textbf{Prioridad Regional 10 de Investigación en Salud de Ayacucho 2025--2030}, aprobada por la Dirección Regional de Salud (DIRESA) de Ayacucho \cite{diresa2025}: \textit{``Organización, gestión, calidad y accesibilidad de los servicios de salud''}, cuyo objetivo estratégico regional apunta a generar evidencia empírica directa para transformar los modelos de gestión y dignificar la atención en los centros de salud de la provincia de Huamanga. En este escenario, el Centro de Salud Belén (Categoría I-3, Microred Huamanga) representa un punto neurálgico asistencial que atiende a una población adscrita de 18,450 habitantes, donde el 68\% se encuentra en situación de pobreza o extrema pobreza según el SISFOH y más del 94\% depende del Seguro Integral de Salud (SIS), coexistiendo un 72\% de población con bilingüismo activo quechua chanka-castellano.
\end{itemize}

\paragraph{2. Importancia y relevancia social del proyecto:}
El estudio aborda un problema crítico y sensible: la distancia existente entre la oferta de servicios programáticos y la vivencia sentida del usuario andino frente al sistema de salud. Al develar los significados, temores, expectativas y vivencias de los usuarios del Centro de Salud Belén, se genera conocimiento de primera fuente para corregir la desatención a la dignidad del paciente, visibilizando la carga emocional y física que soportan las madres de familia en CRED, las gestantes y los adultos mayores durante las horas de espera en las filas de la madrugada.

\paragraph{3. Finalidad e impacto práctico:}
La investigación tiene una clara orientación transformadora: sus hallazgos permitirán a la jefatura del Centro de Salud Belén, a la Dirección de la Red de Salud Huamanga y al colectivo profesional de enfermería diseñar planes de mejora continua, protocolos de acogida y triaje humanizado con pertinencia lingüística intercultural (en quechua y español), y estrategias para desarticular los cuellos de botella en el otorgamiento de citas y entrega de medicamentos.

\paragraph{4. Valor teórico y tipología de la investigación:}
Conforme a la tipología metodológica de la investigación científica, el presente trabajo se clasifica como una \textbf{investigación básica o sustantiva}, con \textbf{nivel descriptivo-interpretativo} y un riguroso \textbf{diseño cualitativo fenomenológico} \cite{sampieri2018, creswell2013}. Siguiendo las directrices metodológicas de Roberto Hernández-Sampieri y Christian Mendoza \cite{sampieri2018}, así como los fundamentos epistemológicos de Edmund Husserl \cite{husserl1949}, Max van Manen \cite{vanmanen1990} y John Creswell \cite{creswell2013}, el propósito central del estudio no es construir modelos teóricos abstractos desvinculados (como en la teoría fundamentada) ni narrar secuencias cronológicas biográficas (como en los diseños narrativos), sino explorar, describir y comprender la \textit{esencia de las experiencias compartidas} (\textit{vivencias}) de los usuarios frente al fenómeno del acceso y la calidad asistencial en el Centro de Salud Belén. Dicho abordaje aprehende la vivencia en sus cuatro dimensiones existenciales fundamentales: temporalidad (tiempo vivido en la espera y la madrugada), espacialidad (espacio vivido en salas y consultorios), corporalidad (cuerpo vivido, dolor, frío y fatiga física) y relacionalidad (encuentro intersubjetivo, empatía y comunicación en quechua chanka). Su valor teórico radica en articular dialógicamente los modelos de Calidad en Salud de Avedis Donabedian \cite{donabedian1980, donabedian2005}, la Teoría del Acceso de Roy Penchansky \cite{penchansky1981} y la Teoría del Cuidado Humano de Jean Watson \cite{watson2008} con la cosmovisión andina de Huamanga.

\paragraph{5. Valor metodológico:}
En el ámbito metodológico, el proyecto aporta una guía de entrevista en profundidad y una matriz de categorización temática cualitativa validadas por juicio de expertos y ajustadas a los criterios de rigor científico de Guba y Lincoln \cite{gubalincoln1985}: credibilidad, transferibilidad, consistencia y confirmabilidad. Asimismo, aplica la técnica de triangulación metodológica y de actores \cite{okuda2005}, sentando un precedente técnico replicable para futuros estudios en la Facultad de Ciencias de la Salud de la UNSCH.

\paragraph{6. Justificación institucional y bioética:}
El protocolo cumple con los estándares epistemológicos promovidos por la cátedra de EN 486 (fundamentados en la concepción científica y rigurosa de Mario Bunge \cite{bunge2014}), respetando de forma irrestricta las normas bioéticas de la Declaración de Helsinki \cite{helsinki2013} y las pautas internacionales del CIOMS \cite{cioms2016}, garantizando la autonomía mediante el consentimiento informado y el estricto anonimato de los participantes.

\newpage
